%% data_integration_report.tex — LIST EuMINe DataBridge 2026
%% Compile: cd reports/data_integration && pdflatex data_integration_report.tex (twice)

\documentclass[10pt, a4paper]{article}
\usepackage{eumine_hackathon}

\title{%
  \textbf{Data Integration Report}\\[4pt]
  {\large EuMINe DataBridge Hackathon \hackathonyear}
}
\author{
  LIST\_EuMINe \\
  \small Luxembourg Institute of Science and Technology (LIST) \\
  \small euminecost@gmail.com
}
\date{June 2026}

\begin{document}
\maketitle
\thispagestyle{fancy}

\begin{abstract}
We integrate three public DFT databases — Materials Project (MP),
JARVIS-DFT, and OQMD — with the EuMINe Bridge Dataset to produce
a harmonised 850-structure training corpus for predicting formation
energy per atom ($E_f$) and electronic band gap ($E_g$).
MP provides primary labels for all 850 Bridge Dataset structures;
JARVIS-DFT supplies 603 matched entries used both for cross-validation
of MP values and to augment ALIGNN pretraining.
Systematic inter-database offsets are quantified (EF offset
$-0.069$~eV/atom, BG offset $+0.099$~eV) and corrected with fitted
linear transformations ($R^2 = 0.961$ and $0.861$ respectively).
\end{abstract}

\tableofcontents
\vspace{1em}

\section{Data Sources Used}

\begin{table}[h!]
\centering
\caption{Data sources integrated into the training dataset.}
\label{tab:sources}
\begin{tabular}{@{}lllll@{}}
\toprule
Source & Version / date & Entries used & Properties & DFT functional \\
\midrule
Bridge Dataset    & 2026 (hackathon) & 850 (train+val) & $E_f$, $E_g$ & PBE/PBE+U \\
Materials Project & v2023.11.1       & 850 (full match) & $E_f$, $E_g$ & PBE (GGA/GGA+U) \\
JARVIS-DFT dft\_3d & 2024-03         & 603 (matched)    & $E_f$, $E_g$ & OptB88vdW \\
OQMD              & v1.5 (REST API)  & 2090 (queried)   & $E_f$        & PBE \\
\bottomrule
\end{tabular}
\end{table}

\subsection{Bridge Dataset}

The Bridge Dataset is the primary source of ground-truth labels,
consisting of 700~training, 150~validation, and 150~test structures
(CIF files), each with formation energy and band gap values sourced
by the hackathon organisers.
Labels were used directly without modification; the CSV
columns \texttt{formation\_energy\_per\_atom\_label} and
\texttt{band\_gap\_label} are taken as the training targets.
Both MP and JARVIS reference values are also provided in the
CSV for the full training set, enabling direct bias analysis
without additional API queries.

\subsection{Materials Project}

All 850 Bridge Dataset structures (train+val) were matched against
Materials Project using the structure \texttt{material\_id} provided
in the CSV.
Data was downloaded via \texttt{mp-api} v0.40+ querying
\texttt{formation\_energy\_per\_atom} and
\texttt{band\_gap} from the \texttt{thermo} and
\texttt{electronic\_structure} endpoints.
MP values served as the primary external reference because the
Bridge Dataset labels are on the same GGA/GGA+U scale.

\subsection{JARVIS-DFT}

The JARVIS-DFT \texttt{dft\_3d} dataset was downloaded via
\texttt{jarvis-tools} and matched to the Bridge Dataset using
the JARVIS ID (\texttt{jid}) column already provided in the
training CSV.
Of the 700 training and 150 validation structures, 603 have
matching JARVIS entries.
JARVIS values use the OptB88vdW functional, which systematically
differs from MP's PBE functional in both $E_f$ and $E_g$
(see Section~\ref{sec:harmonisation}).

Key JARVIS fields used:
\texttt{formation\_energy\_peratom} and \texttt{optb88vdw\_bandgap}.

\subsection{OQMD}

The Open Quantum Materials Database was queried via the OQMD
REST API (\texttt{oqmd.org/optimade}) for structures matching
the Bridge Dataset chemical formulas.
OQMD uses PBE, placing it on a scale close to MP.
OQMD $E_f$ values served as a tertiary fallback in the
harmonisation priority chain (below Bridge label and MP)
but were not used as training targets in the final model.

\section{Harmonisation Strategy}
\label{sec:harmonisation}

\subsection{Structure Matching}

Bridge Dataset $\leftrightarrow$ JARVIS matching uses the JARVIS ID
(\texttt{jid}) column included in the hackathon training CSV,
providing unambiguous one-to-one alignment.
Bridge Dataset $\leftrightarrow$ MP matching uses the MP material ID.
OQMD matching uses reduced formula lookup with a $\pm 2$ site
tolerance.

\subsection{Quantifying Inter-Database Bias}

Using the 603 structures with both MP and JARVIS values, we
computed the systematic offsets:

\begin{table}[h!]
\centering
\caption{Systematic offsets between MP and JARVIS labels.}
\label{tab:offsets}
\begin{tabular}{@{}lccc@{}}
\toprule
Property & Mean offset (MP$-$JAR) & MAE & $R^2$ \\
\midrule
$E_f$ (eV/atom) & $-0.069$ & $0.144$ & $0.961$ \\
$E_g$ (eV)      & $+0.099$ & $0.218$ & $0.861$ \\
\bottomrule
\end{tabular}
\end{table}

The high $E_f$ correlation ($R^2 = 0.961$) indicates that the
systematic offset is well-captured by a linear model.
The lower $E_g$ correlation ($R^2 = 0.861$) is expected: PBE
systematically underestimates band gaps while OptB88vdW partially
corrects this via van der Waals corrections, producing non-uniform
deviations especially for wide-gap insulators.

\subsection{Linear Bias Correction}

A linear correction was fitted by ordinary least squares on the
603 overlapping training pairs:

\begin{align}
E_f^{\text{MP}} &= 1.038 \times E_f^{\text{JARVIS}} - 0.028
    \quad (R^2 = 0.961) \label{eq:ef_corr} \\
E_g^{\text{MP}} &= 0.971 \times E_g^{\text{JARVIS}} + 0.122
    \quad (R^2 = 0.861) \label{eq:bg_corr}
\end{align}

The correction intentionally uses only the Bridge Dataset training
overlap (not test or external data) to prevent leakage.
Parameters are saved to \texttt{data/processed/harmonizer\_params.json}.

\subsection{Label Assignment}

The following priority chain is applied when constructing training
targets (implemented in \texttt{DatabaseHarmonizer.harmonize\_dataframe}):

\begin{enumerate}
  \item Bridge Dataset label — primary, most reliable
  \item MP value — same functional as Bridge label
  \item Corrected JARVIS value (Eq.~\ref{eq:ef_corr}--\ref{eq:bg_corr})
  \item OQMD value ($E_f$ only, PBE scale $\approx$ MP)
  \item AFLOW value (least trusted)
\end{enumerate}

For the final training corpus, all 850 structures have Bridge Dataset
labels; external database values are used as auxiliary inputs
(features) to the ALIGNN fine-tuning targets, not as primary labels.

\section{Data Card}

\begin{table}[h!]
\centering
\caption{Dataset statistics.}
\label{tab:datacard}
\begin{tabular}{@{}ll@{}}
\toprule
Property & Value \\
\midrule
Total entries (train / val / test)   & 700 / 150 / 150 \\
Final training corpus (train + val)  & 850 \\
External DB entries queried          & 2090 (JARVIS), 850 (MP) \\
JARVIS entries matched to Bridge     & 603 / 850 (71\%) \\
Unique reduced formulas              & 694 \\
Unique space groups                  & 103 \\
$E_f$ range (train)                  & $-4.47$ to $+2.36$ eV/atom \\
$E_f$ mean $\pm$ std (train)         & $-1.17 \pm 1.14$ eV/atom \\
$E_g$ mean $\pm$ std (train)         & $0.93 \pm 1.58$ eV \\
Metals ($E_g = 0$, train)            & 428 / 700 (61\%) \\
Semiconductors ($0 < E_g < 3$, train) & 184 / 700 (26\%) \\
Wide-gap ($E_g \geq 3$, train)       & 88 / 700 (13\%) \\
Splitting strategy                   & Stratified by $E_g$ category \\
Median $|\Delta E_f|$ MP–JARVIS      & 0.095 eV/atom \\
Median $|\Delta E_g|$ MP–JARVIS      & 0.140 eV \\
\bottomrule
\end{tabular}
\end{table}

\section{Limitations and Failure Cases}

\begin{itemize}
  \item \textbf{JARVIS coverage gap:} 29\% of Bridge Dataset structures
    (247/850) lack a JARVIS match, preventing bias correction for those
    materials. These are predominantly complex oxides and rare-earth
    compounds that appear sparsely in the JARVIS \texttt{dft\_3d} dataset.
  \item \textbf{DFT band gap underestimation:} Both MP (PBE) and
    JARVIS (OptB88vdW) underestimate $E_g$ for wide-gap insulators.
    This propagates to the training labels for the wide-gap subset
    ($E_g \geq 3$~eV, 13\% of training set) and contributes to the
    elevated BG MAE for this class (0.412~eV on the validation set).
  \item \textbf{Class imbalance:} The training set is dominated by
    metals (61\%), which are easier to predict ($E_g = 0$ trivially).
    Performance estimates are more uncertain for the 13\% wide-gap class
    due to limited training examples.
  \item \textbf{OQMD formula matching:} OQMD matches by formula rather
    than structure identity; for polymorphic materials (multiple
    structures with the same formula), the lowest-energy polymorph
    was selected. This may introduce structure-label mismatches for
    approximately 5--10\% of OQMD-sourced entries.
  \item \textbf{Linear correction simplicity:} The linear harmonisation
    (Eqs.~\ref{eq:ef_corr}--\ref{eq:bg_corr}) cannot capture
    per-element or per-crystal-system biases; more complex corrections
    (e.g.\ per-element linear models) may improve data quality for
    future work.
\end{itemize}

\section*{References}

\begin{enumerate}
  \item Jain et al., \textit{APL Mater.} \textbf{1}, 011002 (2013) —
    Materials Project.
  \item Choudhary et al., \textit{npj Comput. Mater.} \textbf{6}, 173 (2020)
    — JARVIS-DFT.
  \item Saal et al., \textit{JOM} \textbf{65}, 1501 (2013) — OQMD.
  \item Klimeš et al., \textit{J.\ Phys.: Condens.\ Matter} \textbf{22},
    022201 (2010) — OptB88vdW functional.
\end{enumerate}

\end{document}
